\section{Sua parte: base de dados e variáveis}
\textbf{Mensagem principal:} usamos duas fontes oficiais, ligadas pelo código do município: o SIM (homicídios) e o Censo 2022 (características dos municípios).

\subsection{O que enfatizar}
\begin{itemize}
\item \textbf{SIM/DATASUS:} registro oficial de óbitos; contamos agressão (CID-10 X85--Y09) e intervenção legal (Y35--Y36), pelo município onde a vítima morava, somando 2021 a 2023.
\item \textbf{Censo 2022:} retrato de toda a população; renda e trabalho vêm da amostra do Censo (parte dos domicílios), as demais variáveis do questionário completo.
\item \textbf{Variáveis:} a resposta é quantitativa contínua; 10 quantitativas e 3 qualitativas. O estado entra como 26 variáveis indicadoras, com SP como referência.
\item Amazônia Legal: região definida em lei (9 estados). Faixa de fronteira: 150 km ao longo da fronteira terrestre.
\end{itemize}

\subsection{Perguntas prováveis}
\begin{description}[leftmargin=0pt]
\item[Por que somar três anos?] Para estabilizar a taxa: um ano só tem poucos casos em muitos municípios.
\item[Por que município de residência e não de ocorrência?] Porque o denominador é a população residente; assim taxa e população falam das mesmas pessoas.
\item[Por que o Gini não está?] O Gini municipal de 2022 não foi publicado; usamos a razão renda média / mediana como medida de desigualdade.
\end{description}

\subsection{Transição}
Ao terminar as variáveis: ``Agora o Douglas mostra o que os dados revelam antes de qualquer modelo.''
